% ----------------------------------------------------------------------
%  2.2.1. Основные типы дислокаций
% ----------------------------------------------------------------------
\begin{orientir}
Дислокации~— главный «инструмент» пластической деформации кристаллов.
В разделе~\ref{sec:disl} мы разберём: зачем понадобилось понятие
дислокации и что это такое (2.2.1); как дислокации движутся (2.2.2);
какие напряжения создают, какой энергией обладают и какие силы на них
действуют (2.2.3); откуда берутся и как размножаются (2.2.4); какие
группы образуют (2.2.5); как устроены в реальных решётках ГЦК, ОЦК, ГПУ
и как реагируют друг с другом (2.2.6--2.2.7).
\end{orientir}

\subsection{Основные типы дислокаций}\label{ssec:disl-types}

\subsubsection{Загадка прочности: теоретическая прочность на сдвиг}

Пластически деформированный металл покрыт \term{полосами скольжения}:
одна часть кристалла сдвинулась относительно другой по определённой
кристаллографической плоскости. Естественно предположить, что сдвиг
происходит \emph{одновременно} по всей плоскости, как сдвигается одна
половина колоды карт относительно другой. Оценим, какое напряжение для
этого нужно. Эту задачу решил Я.\,И.~Френкель.

\begin{figure}[ht]
\centering
\begin{subfigure}[b]{0.4\textwidth}\centering
\begin{tikzpicture}[scale=0.75]
  \foreach \i in {0,...,5}{\node[atom=6mm] at (\i,0) {};}
  \foreach \i in {0,...,5}{\node[atomB=6mm] at (\i+0.35,1.1) {};}
  \draw[dashed, fmGray] (-0.5,0.55) -- (5.8,0.55);
  \draw[arr, fmRed] (2.2,1.75) -- (3.6,1.75) node[right, lbl]{$\tau$};
  \draw[arr, fmRed] (2.9,-0.65) -- (1.5,-0.65) node[left, lbl]{$\tau$};
  \draw[dimline] (0,1.55) -- node[above, slbl]{$x$} (0.35,1.55);
  \draw[dimline] (0,-0.5) -- node[below, slbl]{$b$} (1,-0.5);
  \draw[dimline] (-0.55,0) -- node[left, slbl]{$a$} (-0.55,1.1);
  \node[slbl] at (6.6,1.1) {I}; \node[slbl] at (6.3,0) {II};
\end{tikzpicture}
\caption{}
\end{subfigure}\hfill
\begin{subfigure}[b]{0.56\textwidth}\centering
\begin{tikzpicture}
\begin{axis}[width=\linewidth, height=4.6cm, xmin=0, xmax=1.05, ymin=-1.25, ymax=1.5,
  axis lines=middle, xtick={0.25,0.5,1}, xticklabels={$b/4$,$b/2$,$b$},
  ytick={1}, yticklabels={$\tau_{\text{теор}}$}, xlabel={$x$}, ylabel={$\tau$},
  every axis x label/.style={at={(ticklabel* cs:1)}, anchor=west},
  every axis y label/.style={at={(ticklabel* cs:1)}, anchor=south}]
  \addplot[fmBlue, very thick, domain=0:1, samples=100] {sin(deg(2*pi*x))};
  \addplot[fmRed, dashed, thick, domain=0:0.22] {2*pi*x};
  \node[fmRed, font=\footnotesize, anchor=west] at (axis cs:0.2,1.33) {закон Гука $\tau=Gx/a$};
\end{axis}
\end{tikzpicture}
\caption{}
\end{subfigure}
\caption{Модель Френкеля: \textit{а}~— атомный ряд I сдвинут
относительно ряда II на $x$ как целое; \textit{б}~— напряжение,
необходимое для сдвига, периодично с периодом $b$ и при малых $x$
подчиняется закону Гука}
\label{fig:frenkel}
\end{figure}

\begin{shagi}
\item Сдвиг на период $b$ возвращает кристалл в эквивалентное состояние,
поэтому напряжение сдвига~— периодическая функция $x$ с периодом $b$.
Простейшее приближение~— синусоида (рис.~\ref{fig:frenkel}):
\[
\tau=\tau_{\text{теор}}\sin\frac{2\pi x}{b}.
\]
\item При малых смещениях выполняется закон Гука: $\tau=G\,x/a$, где
$G$~— модуль сдвига, $a$~— расстояние между атомными плоскостями.
\item При малых $x$ синус равен аргументу: $\tau_{\text{теор}}\cdot2\pi
x/b=Gx/a$.
\end{shagi}
\begin{keyeq}\label{eq:tau-theor}
\tau_{\text{теор}}=\frac{G}{2\pi}\cdot\frac{b}{a}\approx\frac{G}{2\pi}\approx\frac G6 .
\end{keyeq}

Более точные модели межатомного взаимодействия (силы убывают быстрее
синусоиды; возможны дополнительные минимумы энергии) снижают оценку:
Маккензи получил для плотноупакованных решёток $G/30$. Разумный диапазон
\[
\frac G{30}<\tau_{\text{теор}}<\frac G5,\qquad\text{в среднем}\quad\tau_{\text{теор}}\sim\frac G{15}.
\]

\begin{glavnoe}
Реальные кристаллы начинают пластически деформироваться при напряжениях
в тысячи и миллионы раз меньших: $\tau\approx(10^{-3}\text{--}10^{-8})\,G$.
Например, для меди $G\approx48$~ГПа, $\tau_{\text{теор}}\sim3$~ГПа, а
монокристалл чистой меди течёт уже при $\sim1$~МПа. Только почти
совершенные \emph{нитевидные кристаллы} («усы») выдерживают напряжения
$\sim G/15$, как и предсказывает теория.
\end{glavnoe}

Вывод: в обычных кристаллах сдвиг идёт \emph{не одновременно} по всей
плоскости. Он начинается в одном месте и постепенно распространяется~—
так, что в каждый момент смещаются лишь немногие атомы. Граница между
сдвинутой и ещё не сдвинутой частями плоскости и есть
\term{дислокация}. Эту идею высказали в 1934~г. Тейлор, Орован и Поляни.

\begin{smysl}[Аналогия с ковром]
Тяжёлый ковёр трудно сдвинуть, потянув за край: придётся преодолеть
трение по всей площади сразу. Но если сделать на ковре складку и
прогнать её от одного края к другому, ковёр сместится на ширину складки,
а усилие понадобится маленькое: в каждый момент двигается только
складка. Дислокация~— такая «складка» в кристалле. Похоже ползёт и
гусеница: она не сдвигает всё тело сразу, а прогоняет по нему волну.
\end{smysl}

\subsubsection{Краевая дислокация}

Представим, что сдвиг на одно межатомное расстояние прошёл не по всей
плоскости скольжения, а лишь по её части (рис.~\ref{fig:edge-block}).
На участке $ABCD$ плоскости скольжения верхняя часть кристалла уже
сдвинулась на вектор $\vb b$, на остальной части~— ещё нет. Граница
участка сдвига~— линия $AB$~— и есть дислокация.

\begin{figure}[ht]
\centering
\begin{subfigure}[b]{0.48\textwidth}\centering
\begin{tikzpicture}[cubeproj, scale=0.95, line join=round]
  \def\L{4}\def\D{3}\def\H{3}\def\h{1.5}\def\bb{0.5}\def\xd{2}
  % правая грань
  \filldraw[fill=fmBlue!22, draw=atomLine, line width=0.6pt] (\L,0,0)--(\L,\D,0)--(\L,\D,\H)--(\L,0,\H)--cycle;
  % верхняя грань
  \filldraw[fill=fmBlue!8, draw=atomLine, line width=0.6pt] (\bb,0,\H)--(\L,0,\H)--(\L,\D,\H)--(\bb,\D,\H)--cycle;
  % уступ
  \filldraw[fill=fmBlue!8, draw=atomLine, line width=0.6pt] (0,0,\h)--(\bb,0,\h)--(\bb,\D,\h)--(0,\D,\h)--cycle;
  % передняя грань
  \filldraw[fill=fmBlue!14, draw=atomLine, line width=0.6pt] (0,0,0)--(\L,0,0)--(\L,0,\H)--(\bb,0,\H)--(\bb,0,\h)--(0,0,\h)--cycle;
  % область сдвига (внутри)
  \fill[fmAmber, opacity=0.35] (\bb,0,\h)--(\xd,0,\h)--(\xd,\D,\h)--(\bb,\D,\h)--cycle;
  \draw[fmAmber, dashed] (\xd,0,\h)--(\xd,\D,\h)--(\bb,\D,\h);
  \draw[disl, line width=1.8pt] (\xd,0,\h) -- (\xd,\D,\h);
  \node[slbl, below] at (\xd,0,\h) {$A$};
  \node[slbl, above right] at (\xd,\D,\h) {$B$};
  \node[slbl, above left] at (\bb,\D,\h) {$C$};
  \node[slbl, below left] at (0,0,\h) {$D$};
  \draw[bvec] (0.0,0,\h+0.55) -- (\bb+0.02,0,\h+0.55) node[right, slbl, fmRed]{$\vb b$};
  \draw[arr, fmGray] (1.8,0,\H+0.35) -- (3.0,0,\H+0.35) node[right, slbl]{$\tau$};
\end{tikzpicture}
\caption{}\label{fig:edge-block-a}
\end{subfigure}\hfill
\begin{subfigure}[b]{0.48\textwidth}\centering
\begin{tikzpicture}[scale=0.5]
  \begin{scope}\clip (-5.3,-3.6) rectangle (5.3,3.6);
  \EdgeDislLattice{6}{4}{4}\end{scope}
  \dislT{(0,-0.05)}{0}
  \draw[fmGray, dashed] (-5.4,0) -- (5.4,0);
  \node[slbl, fill=white, inner sep=1pt] at (3.9,0) {плоскость скольжения};
\end{tikzpicture}
\caption{}\label{fig:edge-block-b}
\end{subfigure}
\caption{Краевая дислокация. \textit{а}~— сдвиг прошёл только по
участку $ABCD$ плоскости скольжения; граница участка $AB$~— линия
дислокации; вектор сдвига $\vb b$ перпендикулярен $AB$.
\textit{б}~— сечение, перпендикулярное линии дислокации: над плоскостью
скольжения на одну атомную плоскость больше, чем под ней; оранжевым
выделена «лишняя» полуплоскость (экстраплоскость). Положения атомов
рассчитаны по теории упругости}
\label{fig:edge-block}
\end{figure}

Что происходит на атомном уровне, видно из сечения
(рис.~\ref{fig:edge-block},\,б). После сдвига $n$ вертикальных атомных
плоскостей над плоскостью скольжения оказываются напротив $n-1$
плоскостей под ней. Одна плоскость в верхней половине не имеет
продолжения внизу~— это \term{экстраплоскость}.

\begin{opredelenie}
\term{Краевая дислокация}~— линейная область искажённой решётки вдоль
края экстраплоскости. Её линия \emph{перпендикулярна} вектору сдвига:
$\vb b\perp\vb t$, где $\vb t$~— единичный вектор вдоль линии дислокации.
\end{opredelenie}

\begin{itemize}
\item Над краем экстраплоскости межатомные расстояния меньше нормальных
(\emph{сжатие}), под ним~— больше (\emph{растяжение}). Атом на самой кромке
экстраплоскости имеет меньше соседей, чем в совершенном кристалле.
\item Сильные искажения сосредоточены в трубке радиусом в несколько
межатомных расстояний вокруг края экстраплоскости~— это \term{ядро
дислокации}. Вдоль линии дислокации ядро тянется на всю её длину.
Точное расположение атомов в ядре рассчитывают методами компьютерного
моделирования.
\item Если экстраплоскость вставлена сверху, дислокацию обозначают
$\perp$ и условно называют \emph{положительной}; если снизу~— $\top$,
\emph{отрицательной}. «Ножка» значка направлена в сторону лишнего
материала. Знак~— дело соглашения: перевернув кристалл, мы превратим
положительную дислокацию в отрицательную.
\end{itemize}

\subsubsection{Винтовая дислокация}

Сделаем теперь надрез $ABCD$, дошедший только до середины кристалла, и
сдвинем одну его сторону относительно другой на межатомное расстояние
\emph{вдоль} края надреза (рис.~\ref{fig:screw-block}).

\begin{figure}[ht]
\centering
\begin{subfigure}[b]{0.48\textwidth}\centering
\begin{tikzpicture}[cubeproj, scale=0.95, line join=round]
  \def\L{4}\def\D{3}\def\H{3}\def\bb{0.5}\def\xc{2}\def\yc{1.5}
  % правая грань (передняя часть опущена на b)
  \filldraw[fill=fmBlue!22, draw=atomLine, line width=0.6pt]
     (\L,0,-\bb)--(\L,\yc,0)--(\L,\D,0)--(\L,\D,\H)--(\L,\yc,\H)--(\L,0,\H-\bb)--cycle;
  % верх: левая часть и задняя правая
  \filldraw[fill=fmBlue!8, draw=atomLine, line width=0.6pt]
     (0,0,\H)--(\xc,0,\H)--(\xc,\yc,\H)--(\L,\yc,\H)--(\L,\D,\H)--(0,\D,\H)--cycle;
  % пандус (передняя правая часть верха)
  \filldraw[fill=fmBlue!12, draw=atomLine, line width=0.6pt]
     (\xc,0,\H-\bb)--(\L,0,\H-\bb)--(\L,\yc,\H)--(\xc,\yc,\H)--cycle;
  % ступенька
  \filldraw[fill=fmAmber!45, draw=atomLine, line width=0.6pt]
     (\xc,0,\H)--(\xc,0,\H-\bb)--(\xc,\yc,\H)--cycle;
  % передняя грань
  \filldraw[fill=fmBlue!14, draw=atomLine, line width=0.6pt] (0,0,0)--(\xc,0,0)--(\xc,0,\H)--(0,0,\H)--cycle;
  \filldraw[fill=fmBlue!14, draw=atomLine, line width=0.6pt] (\xc,0,-\bb)--(\L,0,-\bb)--(\L,0,\H-\bb)--(\xc,0,\H-\bb)--cycle;
  % плоскость надреза и линия дислокации
  \fill[fmAmber, opacity=0.25] (\xc,0,0)--(\xc,\yc,0)--(\xc,\yc,\H)--(\xc,0,\H)--cycle;
  \draw[disl, line width=1.8pt, dashed] (\xc,\yc,\H) -- (\xc,\yc,0);
  \fill[fmBlue] (\xc,\yc,\H) circle (1.6pt);
  \node[slbl, above left] at (\xc,0,\H) {$A$};
  \node[slbl, above] at (\xc,\yc,\H) {$B$};
  \node[slbl, right] at (\xc,\yc,0.05) {$C$};
  \node[slbl, below left] at (\xc,0,0) {$D$};
  \draw[bvec] (\xc+0.25,0,\H+0.05) -- (\xc+0.25,0,\H-\bb-0.02);
  \node[slbl, fmRed, right] at (\xc+0.3,0,\H+0.15) {$\vb b$};
\end{tikzpicture}
\caption{}
\end{subfigure}\hfill
\begin{subfigure}[b]{0.48\textwidth}\centering
\begin{tikzpicture}[scale=0.42]
  \ScrewDislPlan{5}{6}
  \draw[disl, line width=1.2pt] (0,-0.6) -- (0,6.6);
  \node[slbl, fmBlue, above] at (0,6.6) {линия дислокации};
  \node[slbl, align=center] at (-3,-1.4) {сдвиг прошёл};
  \node[slbl, align=center] at (3,-1.4) {сдвига нет};
\end{tikzpicture}
\caption{}
\end{subfigure}
\caption{Винтовая дислокация. \textit{а}~— надрез $ABCD$ и сдвиг на
$\vb b$ параллельно краю надреза; линия дислокации $BC$ параллельна
$\vb b$; на верхней грани~— ступенька, сходящая на нет в точке $B$.
\textit{б}~— атомы двух слоёв по обе стороны плоскости сдвига, вид на
эту плоскость: тёмные~— под ней, светлые~— над ней. Слева (сдвиг прошёл)
и справа (сдвига нет) атомы двух слоёв снова друг над другом; вблизи
линии дислокации они смещены, и ряды образуют винтовую поверхность}
\label{fig:screw-block}
\end{figure}

На верхней грани появилась ступенька, которая не проходит через весь
кристалл, а сходит на нет в точке $B$. Сдвинуться целиком на период
правая часть не может: надрез дошёл лишь до середины. Поэтому смещение
убывает от $A$ к $B$, и горизонтальные атомные плоскости изгибаются. В
итоге кристалл, состоявший из параллельных атомных слоёв, превращается
в \emph{одну} атомную плоскость, закрученную винтовой лестницей
(геликоидом). Обходя вокруг линии $BC$, мы поднимаемся (или опускаемся)
на одну ступеньку винта~— межатомное расстояние.

\begin{opredelenie}
\term{Винтовая дислокация}~— линейный дефект, вокруг линии которого
атомные плоскости закручены винтом. Её линия \emph{параллельна} вектору
сдвига: $\vb b\parallel\vb t$.
\end{opredelenie}

Винтовая дислокация бывает \emph{правой} ($\vb b\uparrow\uparrow\vb t$)
и \emph{левой} ($\vb b\uparrow\downarrow\vb t$)~— как правая и левая
резьба. В отличие от краевой, правую винтовую дислокацию нельзя
превратить в левую, просто перевернув кристалл: правый винт при
переворачивании остаётся правым.

\subsubsection{Вектор Бюргерса}

Тип и «мощность» дислокации характеризует \term{вектор Бюргерса}
$\vb b$. Его определяют с помощью \term{контура Бюргерса}~— замкнутого
пути из шагов «от атома к атому» (рис.~\ref{fig:burgers}).

\begin{figure}[ht]
\centering
\begin{subfigure}[b]{0.48\textwidth}\centering
\begin{tikzpicture}[scale=0.5]
  \begin{scope}\clip (-4.6,-3.6) rectangle (4.6,3.6);
  \EdgeDislLattice{5}{4}{4}\end{scope}
  \dislT{(0,-0.05)}{0}
  \draw[fmRed, line width=1.1pt, line join=round,
    decoration={markings, mark=at position 0.12 with {\arrow{Stealth[length=2.2mm]}},
      mark=at position 0.4 with {\arrow{Stealth[length=2.2mm]}},
      mark=at position 0.66 with {\arrow{Stealth[length=2.2mm]}},
      mark=at position 0.9 with {\arrow{Stealth[length=2.2mm]}}}, postaction={decorate}]
    (L-3_-3) -- (L-3_-2) -- (L-3_-1) -- (U-3_0) -- (U-3_1) -- (U-3_2)
    -- (U-2_2) -- (U-1_2) -- (U0_2) -- (U1_2) -- (U2_2) -- (U3_2)
    -- (U3_1) -- (U3_0) -- (L2_-1) -- (L2_-2) -- (L2_-3)
    -- (L1_-3) -- (L0_-3) -- (L-1_-3) -- (L-2_-3) -- (L-3_-3);
  \node[fmRed, slbl, below left] at (L-3_-3) {$S{=}F$};
\end{tikzpicture}
\caption{реальный кристалл}
\end{subfigure}\hfill
\begin{subfigure}[b]{0.48\textwidth}\centering
\begin{tikzpicture}[scale=0.5]
  \foreach \i in {-4,...,4}{\foreach \j in {-3,...,3}{
     \draw[atomLine!45, line width=0.5pt] (\i-0.5,\j) -- (\i+0.5,\j) (\i,\j-0.5) -- (\i,\j+0.5);}}
  \foreach \i in {-4,...,4}{\foreach \j in {-3,...,3}{
     \fill[atomA, draw=atomLine, line width=0.4pt] (\i,\j) circle (0.2);}}
  \draw[fmRed, line width=1.1pt, line join=round,
    decoration={markings, mark=at position 0.12 with {\arrow{Stealth[length=2.2mm]}},
      mark=at position 0.4 with {\arrow{Stealth[length=2.2mm]}},
      mark=at position 0.66 with {\arrow{Stealth[length=2.2mm]}},
      mark=at position 0.9 with {\arrow{Stealth[length=2.2mm]}}}, postaction={decorate}]
    (-3,-3) -- (-3,2) -- (3,2) -- (3,-3) -- (-2,-3);
  \draw[bvec, line width=1.6pt] (-2,-3) -- (-3,-3);
  \node[slbl, below] at (-3,-3.2) {$S$};
  \node[slbl, below] at (-2,-3.2) {$F$};
  \node[slbl, fmRed, above] at (-2.5,-2.95) {$\vb b$};
\end{tikzpicture}
\caption{совершенный кристалл}
\end{subfigure}
\caption{Контур Бюргерса. \textit{а}~— в реальном кристалле контур
обходит дислокацию по «хорошему» материалу и замыкается: 5~шагов вверх,
6~вправо, 5~вниз, 5~влево. \textit{б}~— тот же контур в совершенном
кристалле не замыкается; замыкающий вектор $\vb b=\overrightarrow{FS}$~—
вектор Бюргерса}
\label{fig:burgers}
\end{figure}

\paragraph{Правило построения (FS/RH)}
\begin{shagi}
\item Выбираем направление линии дислокации $\vb t$ (произвольно; пусть
$\vb t$ направлен за плоскость рисунка).
\item В реальном кристалле строим замкнутый контур, обходя дислокацию по
часовой стрелке (по правому винту относительно $\vb t$). Контур проходит
только через «хороший» материал, где решётка искажена лишь слегка.
\item Повторяем те же шаги (столько же шагов в тех же направлениях) в
совершенном кристалле. Контур не замкнётся: он начнётся в точке $S$
(\emph{start}) и закончится в $F$ (\emph{finish}).
\item Вектор, замыкающий контур от конца к началу, $\overrightarrow{FS}$,~—
вектор Бюргерса $\vb b$ (правило «конец--начало по правому винту»,
FS/RH).
\end{shagi}

Отсюда ещё одно определение: \emph{дислокация~— линейный дефект, который
«размыкает» любой контур Бюргерса, охватывающий её линию}. Контур, не
охватывающий дислокацию, замыкается и в реальном, и в совершенном
кристалле.

\begin{vnimanie}
Знак $\vb b$ зависит от выбора направления $\vb t$: если изменить
направление линии на противоположное, вектор Бюргерса тоже сменит
направление. В литературе нет единого соглашения о знаке, поэтому важно
сразу оговаривать, какое правило используется.
\end{vnimanie}

\begin{podrobnee}[Вектор Бюргерса в сплошной среде]
Через поле упругих смещений $\vb u$ вектор Бюргерса записывается как
интеграл по замкнутому контуру $L$, охватывающему линию дислокации:
\[
\vb b=\oint_L\frac{\partial\vb u}{\partial l}\,\d l .
\]
В сплошной среде дислокацию можно создать так (\emph{конструкция
Вольтерры}): вырезать вдоль будущей линии дислокации тонкий цилиндр,
разрезать тело до этого цилиндра, сместить берега разреза друг
относительно друга на $\vb b$ и склеить. После «склейки» любой контур
вокруг оси не возвращается в исходную точку. Тело с дислокацией
топологически отличается от сплошного шара~— оно \emph{многосвязное}.
\end{podrobnee}

\paragraph{Типы дислокаций через $\vb b$ и $\vb t$}
Теперь типы дислокаций можно определить формально
(табл.~\ref{tab:disl-types}).

\begin{table}[ht]
\centering
\caption{Типы дислокаций}
\label{tab:disl-types}
\small
\begin{tabularx}{\textwidth}{@{}lll>{\raggedright\arraybackslash}X@{}}
\toprule
Тип & Условие & Расположение & Плоскость скольжения \\
\midrule
краевая & $\vb b\cdot\vb t=0$ & $\vb b\perp\vb t$ & единственная (содержит $\vb b$ и $\vb t$) \\
правая винтовая & $\vb b\cdot\vb t=b$ & $\vb b\uparrow\uparrow\vb t$ & любая, содержащая линию \\
левая винтовая & $\vb b\cdot\vb t=-b$ & $\vb b\uparrow\downarrow\vb t$ & любая, содержащая линию \\
смешанная & $0<|\vb b\cdot\vb t|<b$ & под углом & единственная \\
\bottomrule
\end{tabularx}
\end{table}

Дислокации с одинаковыми $\vb t$ и $\vb b$ (или с $\vb t_1=-\vb t_2$ и
$\vb b_1=-\vb b_2$, что то же самое) называют \term{одноимёнными}.

\subsubsection{Основные свойства дислокаций}

Контуры Бюргерса, которые можно совместить друг с другом деформацией и
переносом, не пересекая «плохой» области, называют
\emph{эквивалентными}; они дают один и тот же $\vb b$. Сдвигая контур
вдоль дислокации, получаем главное свойство:

\begin{glavnoe}
\begin{enumerate}[label=\arabic*)]
\item Вектор Бюргерса \vazhno{одинаков вдоль всей линии дислокации},
как бы она ни изгибалась. Меняется только ориентация линии, а с ней и
тип участка (краевой, винтовой, смешанный).
\item Поэтому дислокация \vazhno{не может оборваться внутри кристалла}.
Она либо выходит концами на поверхность, либо замыкается в петлю, либо
разветвляется в \term{узле} на другие дислокации.
\item В узле выполняется аналог правила Кирхгофа: если все линии
направить от узла, то $\sum_i\vb b_i=0$.
\end{enumerate}
\end{glavnoe}

\begin{masshtab}[дислокация]{-9.6/-9/{ядро $\lesssim1$ нм}, -8/-5/{расстояние между дислокациями}, -2.5/-1.5/{образец 1 см}}
Ядро дислокации~— «трубка» диаметром в несколько $b$ ($\lesssim1$~нм);
упругое поле убывает как $1/r$ и тянется до соседних дислокаций.
Расстояние между дислокациями $\sim1/\sqrt\rho$: в отожжённом металле
($\rho\sim10^{10}$--$10^{12}$~м$^{-2}$)~— 1--10~мкм, в сильно деформированном
($\rho\sim10^{15}$--$10^{16}$~м$^{-2}$)~— 10--30~нм. При этом линии
макроскопически длинные: плотность~— это длина линий в единице объёма,
и при $\rho=10^{12}$~м$^{-2}$ в 1~см$^3$ металла $10^{12}\cdot10^{-6}=10^6$~м,
т.\,е. тысяча километров дислокаций.
\end{masshtab}

Хорошо это видно на дислокационной \emph{петле}, лежащей в плоскости
скольжения (рис.~\ref{fig:loop-types}). Внутри петли сдвиг прошёл,
снаружи~— нет. Вектор $\vb b$ всюду один, а линия поворачивает. Там, где
линия перпендикулярна $\vb b$, участок краевой (с противоположных сторон
петли~— противоположного знака); где параллельна~— винтовой (правый и
левый); в промежуточных точках~— смешанный.

\begin{figure}[ht]
\centering
\begin{tikzpicture}[scale=0.9]
  \fill[fmBlue!6] (-3.6,-2.4) rectangle (3.6,2.4);
  \draw[fmGray] (-3.6,-2.4) rectangle (3.6,2.4);
  \fill[fmAmber!25] (0,0) circle (1.7);
  \draw[disl, line width=1.6pt, decoration={markings, mark=at position 0.3 with {\arrow{Stealth[length=2.6mm]}}}, postaction={decorate}] (0,0) circle (1.7);
  \node[slbl, align=center] at (0,0) {сдвиг\\прошёл};
  \draw[bvec, line width=1.5pt] (-3.2,1.95) -- (-2.0,1.95) node[right, fmRed, slbl]{$\vb b$};
  \dislT[fmRed]{(1.7,0)}{-90}
  \dislT[fmRed]{(-1.7,0)}{90}
  \node[slbl, right] at (2.1,0) {краевой};
  \node[slbl, left] at (-2.1,0) {краевой};
  \node[slbl, above] at (0,1.75) {винтовой};
  \node[slbl, below] at (0,-1.75) {винтовой};
  \node[slbl, anchor=south west] at (1.25,1.2) {смешанный};
  \node[slbl, anchor=north east] at (-1.25,-1.2) {смешанный};
\end{tikzpicture}
\caption{Дислокационная петля в плоскости скольжения (вид на плоскость
сверху). Вектор Бюргерса один для всей петли; тип участка определяется
углом между линией и $\vb b$}
\label{fig:loop-types}
\end{figure}

\begin{itogi}
\begin{itemize}
\item Теоретическая прочность на сдвиг $\sim G/15$, реальная~— в
$10^3$--$10^8$ раз меньше. Объяснение: сдвиг идёт постепенно, его фронт~—
дислокация.
\item Дислокация~— граница между сдвинутой и несдвинутой частями плоскости
скольжения. Краевая: $\vb b\perp\vb t$, есть экстраплоскость. Винтовая:
$\vb b\parallel\vb t$, плоскости закручены винтом. Смешанная~— всё
остальное.
\item Вектор Бюргерса~— замыкающий вектор контура Бюргерса (правило
FS/RH); он постоянен вдоль линии.
\item Дислокация не обрывается внутри кристалла; в узлах $\sum\vb b_i=0$.
\end{itemize}
\end{itogi}

\begin{proverka}
\item Почему реальная прочность кристаллов намного меньше теоретической?
\item Как по векторам $\vb b$ и $\vb t$ определить тип дислокации?
\item Может ли одна и та же дислокация быть на одном участке краевой, а
на другом винтовой?
\item Почему дислокация не может закончиться внутри кристалла?
\item Три дислокации сходятся в узле. Векторы Бюргерса двух из них (все
линии направлены от узла) $\vb b_1$ и $\vb b_2$. Чему равен $\vb b_3$?
\end{proverka}
